% ===== 原文 PDF pp.32--34（印刷 pp.31--33）=====
% 译者 tr-03；承接 tr-02（p.31 末含 §3 标题与引子），自 p.32 首行直接续写
% [原文 p.31]
% 【接缝】p.31 末句（“…these Abelian vector multiplets correspond to the Cartan generators …”
%  含 $\langle q\rangle=\langle\tilde q\rangle=0$ 与 $\langle\Phi\rangle$）已由 tr-02 在其
%  tr02c.tex 末尾写完并标 `% [边界] 本句已写完`；本块自 p.32 第二句起，不重复该句。
对 $G=\SU(r+1)$，在相差一个规范变换的意义下，该真空期望值形如
\begin{equation}
  \langle\Phi\rangle = \diag\left(a_1,\,a_2,\,\ldots,\,a_{r+1}=-\sum_i a_i\right)\,.
  \tag{3.1}
\end{equation}
非零的 $a_i$ 照例给出有质量的 $W$ 玻色子，$m_W\sim|a_i-a_j|$；通过超势 $W=\tilde Q\Phi Q$
给出有质量的超多重态，$m_H\sim|a_i|$；以及不那么显然的、有质量的磁单极子，
$m_M\sim\frac{1}{g_i^2}|a_i|$，其中 $g_i$ 表示麦克斯韦耦合。我们感兴趣的是把所有这些有质量
场积分掉后得到的、无质量 $\mathrm{U}(1)^r$ 矢量多重态的威尔逊有效场论。\footnote{我们强调，
这个威尔逊有效场论的存在与是否有一个拉格朗日型的 UV 描述无关。不过，初次接触的读者不妨
借助拉格朗日量的例子来理解。}下文在上下文不致产生混淆时，将稍微滥用记号：用
$a_i\ (i=1,2,\ldots,r)$ 同时表示阿贝尔矢量多重态及其最下分量标量。

% [Tutor Note: 这里“最下分量”(bottom component) 指超多重态中 $\theta=\bar\theta=0$ 的分量，
% 即标量 $a_i$；下文 (3.22) 中“$a$ 与 $a^D$ 是对应（对偶）$\cN=2$ 矢量多重态的最下分量”用的是同一说法。
% 原文在同一段里先用 $a_i$ 表示整个多重态、再用它表示标量，这是超场文献的常规滥用，不是笔误。]

该有效场论由 CB 上一般点处的局域拉格朗日量 $\cL_{\mathrm{EFT}}$ 支配。尽管以常数耦合的
麦克斯韦作用量所描述的普通纯阿贝尔规范理论是平凡的，这个 CB 有效场论却因额外的标量场而
具有丰富的动力学——麦克斯韦耦合会随这些标量场以非平凡的方式变化。一般而言，威尔逊有效
作用量是一个非常复杂的对象。然而，$\cN=2$ 超对称对 CB 上的有效拉格朗日量
$\cL_{\mathrm{EFT}}$ 施加了严格的限制，\footnote{在具有 $\cN=2$ 超对称的情形，CB 上的威尔逊
作用量在二阶导数层次上与 1PI 有效作用量一致。进一步的讨论见 [84] 及其中的文献。}且在
$\cN=1$ 超空间中它取如下形式 [1,121,148]
\begin{equation}
  \resizebox{0.82\textwidth}{!}{$\displaystyle
  \cL_{\mathrm{EFT}} \supset \frac{1}{4\pi}\Im\left(
  \int d^4\theta\,\frac{\partial\cF}{\partial\Phi_i}\bar\Phi_i
  + \int d^2\theta\,\frac{1}{2}\frac{\partial^2\cF}{\partial\Phi_i\partial\Phi_j}
  W^i_\alpha W^{j\alpha}\right),\qquad i,j=1,\ldots,r$}\,,
  \tag{3.2}
\end{equation}
其中 $\Phi_i$ 与 $W^i_\alpha$ 分别表示 $\cN=2$ 的 $\mathrm{U}(1)^r$ 矢量多重态 $a_i$ 中的手征
超场与规范微子超场。$\cF$ 是一个（局域的）全纯函数，称为预势；由于超对称，它决定了整个
有效作用量（至二阶导数）。这样的限制把特殊 Kähler 几何与 4d $\cN=2$ CB 理论的动力学紧密
联系起来。为看清这一点，回顾 (3.2) 右端第一项就是所谓的 Kähler 势
$K(\Phi_i)$，而把 VEV (3.1) 提升为

% [Tutor Note: (3.2) 中 $\Phi$ 的指标是下标 $i$、$W$ 的指标是上标 $i$ 加下标 $\alpha$
% （原文合写为 $W^i_\alpha W^{j\alpha}$，第二个因子把 $\alpha$ 上提后缩并）。文字层会把这种
% 上下标关系抽错，此处已按 800 dpi 裁图逐符号核对。]

% [原文 p.32]
缓变的模场之后，其标量部分化为
\begin{equation}
  \Im\left(\frac{\partial^2\cF}{\partial a_i\partial a_j}\,da^i\,d\bar a^{\bar j}\right)
  \tag{3.3}
\end{equation}
其中 $\bar a$ 表示 $a$ 的复共轭。以非线性 sigma 模型的语言，它定义了不等价真空构型空间
（即模空间 $\cM_{\mathrm{CB}}$）上的度规
\begin{equation}
  g_{i\bar j} = \partial_i\partial_{\bar j}K
  = \Im\left(\frac{\partial^2\cF}{\partial a_i\partial a_j}\right)\,,
  \tag{3.4}
\end{equation}
而该空间还是一个 $r$ 复维的刚性特殊 Kähler 流形，因为度规由预势 $\cF$ 确定 [149]。幺正性要求
标量动能项为正，因此预势 $\cF$ 受到约束，使 sigma 模型度规 (3.4) 正定。

% [Tutor Note: (3.4) 只需一步：把 (3.2) 中被 $\int d^4\theta$ 积分的
% $\frac{\partial\cF}{\partial\Phi_i}\bar\Phi_i$ 认作 Kähler 势 $K$，则 $g_{i\bar j}=\partial_i\partial_{\bar j}K$；
% 对 $\bar\Phi_j$ 求导恰好消去 $\bar\Phi$ 因子、只留下全纯-全纯二阶导 $\partial_i\partial_j\cF$。
% 这也解释了为何 (3.4) 左端带 $\bar j$、括号里两个求导都不带横线。]

(3.2) 右端第二项包含 $\mathrm{U}(1)^r$ 规范场的麦克斯韦拉格朗日量，
\begin{equation}
  \cL_{EM} = \frac{1}{16\pi}\left[
  \Im\left(\frac{\partial^2\cF}{\partial a_i\partial a_j}\right)F^i\wedge\star F^j
  + \Re\left(\frac{\partial^2\cF}{\partial a_i\partial a_j}\right)F^i\wedge F^j\right].
  \tag{3.5}
\end{equation}
因此，预势 $\cF$ 也通过
$\tau^{\mathrm{IR}}=\frac{\theta}{2\pi}+\frac{4\pi i}{g^2}$ 决定了依赖于场的复化规范耦合
（矩阵）
\begin{equation}
  \tau^{\mathrm{IR}}_{ij} = \frac{\partial^2\cF}{\partial a_i\partial a_j}\,,
  \tag{3.6}
\end{equation}
其虚部与 sigma 模型度规 (3.4) 一致，而虚部的正定性保证规范动能项是良定义的。

有一点需要说明：到目前为止，我们已经看到 CB 模空间 $\cM_{\mathrm{CB}}$ 的两组坐标，即上文
讨论的矢量多重态标量 $a_i$，以及如定义 2.3 所指出的库仑分支算子
\begin{equation}
  u_i := \Tr(\Phi^{i+1})\,,\qquad i=1,\ldots,r\,,
  \tag{3.7}
\end{equation}
由规范不变算子定义的坐标 $\{u_i\}$ 是物理的、无歧义的。相反，坐标 $\{a_i\}$ 不是规范不变的
（它们在 $\SU(N)$ 的 Weyl 群 $S_N$ 下变换），并且由于 CB 上的电磁对偶还带有进一步的不确定性，
这一点我们很快就会谈到。因此，$\{a_i\}$ 只在 CB 上局域定义，从一个开片到另一个开片要经受
规范变换与对偶变换的等同。尽管如此，这些局域坐标

% [Tutor Note: $u_i=\Tr(\Phi^{i+1})$ 的次数为 $i+1$；$i=1,\ldots,r$ 给出 $r$ 个独立不变量，
% 与 CB 的 $r$ 个复维相配。$i=1$ 时给出 $u=\Tr\Phi^2$（下文 (3.16) 之后对秩一情形直接用
% $u\equiv\Tr(\Phi^2)$），偶次幂；$\Tr\Phi$ 因 $\fsu(N)$ 无迹而恒为零，故求和自 $i+1=2$ 起。]

% [原文 p.33]
$\{a_i\}$ 正是使 $\cN=2$ CB 上的 EFT 描述成为可能、并编码出一个优雅的涌现几何的东西：坐标中
的冗余导致了对该几何的约束。

在 UV 中，预势 $\cF_{\mathrm{UV}}$ 可以从 $\cN=2$ 超杨--米尔斯理论的经典拉格朗日量项 (2.23)
读出，其形式为
\begin{equation}
  \cF_{\mathrm{UV}} \sim \tau_{\mathrm{UV}}\Tr\left(\mathbf V_{\cN=2}^2\right),
  \tag{3.8}
\end{equation}
其中 $\mathbf V_{\cN=2}$ 表示 $\cN=2$ 超空间形式下的 $\cN=2$ 矢量多重态，\footnote{这里我们为
$\cN=2$ 超空间引入另一组超坐标 $\tilde\theta$，而经典拉格朗日量 (2.23) 可简单地写成
\begin{equation}
  \cL_{\mathrm{EFT}} = \int d^2\theta\,d^2\tilde\theta\,\cF(a_i) + \mathrm{h.c.}\,.
  \tag{3.9}
\end{equation}
关于这一形式的更多细节可参见例如 [150]。}它把 $\cN=1$ 手征多重态 $Q$ 与 $\cN=1$ 规范微子
多重态 $W_\alpha$ 紧凑地打包为
$\mathbf V_{\cN=2} := Q + W_\alpha\tilde\theta^\alpha$。在 IR 中，预势 $\cF$ 会更为复杂，它是
矢量多重态标量 $a_i$ 的全纯函数，带有各种可能依赖于动力学标度 $\Lambda$ 与质量参数的修正，
\begin{equation}
  \cF(a_i) \sim \tau^{\mathrm{IR}}_{ij}(a)\,a^i a^j + \ldots
  \tag{3.10}
\end{equation}
其中当能量标度 $E$ 取为 $a$ 时，$\tau^{\mathrm{IR}}_{ij}(a)$ 具有 (2.40) 形式的展开。SW
理论的主要目标，是从 UV 中相应的预势确定 IR 预势；正如 Seiberg 与 Witten 所指出的
[20,21]，在给定某些物理输入之后，它可以由相应的 SW 几何完全确定。\footnote{另一方面，给定
一个一般的 SW 几何（它编码了某个假定的 SCFT 的低能有效场论），要提取相应 SCFT 的信息可能
很困难。关于利用 SW 几何纤维的混合 Hodge 结构的最新尝试，参见 [151]。}

本节其余部分的组织如下。在 3.1 节我们回顾 4d $\cN=2$ 理论中的电磁对偶。在 3.2 节我们介绍
4d $\cN=2$ SCFT 的预势及其在描述 CB 上特殊几何中的作用。在 3.3 节我们解释如何从 SW 曲线
得到量子预势。我们在 3.4 节介绍 Argyres--Douglas（AD）理论及其标度不变的 SW 解。最后，在
3.5 节我们描述可以从 AD 点附近的 CB 有效场论中提取的 CFT 数据。

% [Tutor Note: (3.8) 与 (3.9) 是同一件事的两种写法：在 $\cN=2$ 超空间中预势 $\cF$ 直接充当
% 拉格朗日密度被 $\int d^2\theta\,d^2\tilde\theta$ 积分，而在 $\cN=1$ 超空间中它需经 (3.2) 的
% 两个 $\int d^4\theta$、$\int d^2\theta$ 项才给出同样的作用量；这解释了为什么下一句说 $\cF$
% “决定了整个有效作用量（至二阶导数）”。]

\SUB{3.1}{电磁对偶}

我们已经看到，CB 上的有效场论由 $\mathrm{U}(1)^r$ 矢量多重态的拉格朗日量 (3.2) 局域地支配。
事实证明，这一拉格朗日描述并不唯一。
